\chapter{Introduction}
\label{ch:introduction}

\section{Background of the Study}
\label{sec:background}

Globally, universities are transitioning from traditional, paper-based administrative workflows to digital ecosystems to enhance student access to educational funding. This digitalization is increasingly necessary due to rising educational costs and expanding student populations. In the Philippines, the steady growth of the higher education sector \cite{ref29} has coincided with tuition fee increases across numerous private institutions, driven by inflation and operational adjustments \cite{ref11,ref19}. While financial aid is available through state-sponsored grants \cite{ref21,ref44}, international educational investments \cite{ref45}, and private corporate foundations \cite{ref1,ref40}, the institutions that administer this aid must still collect, verify, and deliberate on large volumes of applicant documents. At Ateneo de Davao University (AdDU), students face similar financial pressures, exacerbated by recent and proposed tuition increases. AdDU's wider financial aid landscape spans internally funded endowments, external corporate partnerships, and state-sponsored grants \cite{ref2,ref8} (Appendix~\ref{app:pipelines}). Within that landscape, the university itself selects scholars through three internal tracks, which together receive approximately 1,326 applicants per cycle (Table~\ref{tab:tracks}). External and government grants fall outside the university's selection process; AdDU only processes their disbursement after receiving each grantor's official beneficiary roster.

% TODO(OPEN-5): currency and period of the income limit; unit of the minimum GIA grant.
% TODO(OPEN-3), TODO(OPEN-9): Jubilee guarantee vs. program minimums; SOP "academic grant" = Jubilee?
\begin{table}[htbp]
  \centering
  \small
  \caption{Internal Scholarship Tracks Governed by ScholarPath AdDU}
  \label{tab:tracks}
  \begin{tabular}{@{}p{2.3cm}p{3.6cm}p{3.6cm}p{4.9cm}@{}}
    \toprule
    Track & Beneficiaries & Grant & Key Rules \\
    \midrule
    Jubilee Scholarship & Certified high school Rank~1 and Rank~2 graduates & Fixed: Rank~1 (Valedictorian) 100\% tuition; Rank~2 (Salutatorian) 50\% tuition & Requires an official ranking certificate; rates are fixed; program-specific admission minimums are applied by each school \\
    Grant-in-Aid (GIA) & Financially challenged students below the university's household income ceiling & Variable, up to 100\% tuition & Requires verified income documents (BIR Form 2316 or sworn statement); slots are highly competitive \\
    Working Scholars & Need-based applicants willing to render service & 100\% tuition, fulfilled through assigned campus work hours instead of cash & Same initial screening as GIA; work obligations are tracked in the Academic Trajectory Path \\
    \bottomrule
  \end{tabular}
\end{table}

Admission to these tracks follows the university's Standard Procedure for Scholarship Applications (SOP) \cite{ref51}, an eight-step process that runs from the application announcement through document submission, verification, interview, deliberation by the School Scholarship Subcommittee, and the release of results by the Office of Admission. Two features of this process create persistent operational strain. First, document collection is single-phase and paper-heavy: every applicant submits a complete physical file, and an application with a missing or illegible document can be rejected outright. Second, final selection belongs to five School Scholarship Subcommittees, those of the School of Nursing (SON), the School of Engineering and Architecture (SEA), the School of Business and Governance (SBG), the School of Education (SOE), and the School of Arts and Sciences (SAS). Each subcommittee applies its own selection criteria, so applicant data must be collated manually for each school.

To address this, this project proposes ScholarPath AdDU, a Data-Feeding Committee Portal for AdDU's internal scholarship admission process. The platform is deliberately not an automated matching tool. It structures, verifies, and presents applicant data to the people who decide, and it records every decision that changes an applicant's status together with the person who made it. The name ScholarPath refers to two structured paths: the Application Lifecycle Path, which carries an applicant from registration to an award, waitlist, or disapproval outcome, and the Academic Trajectory Path, which tracks an awarded scholar's academic standing, service obligations, and graduation. The platform's first core component is a two-stage digital submission into a Document Vault backed by Supabase Storage. Phase~1 Pre-qualification collects the documents needed to screen basic eligibility, and Phase~2 Full Verification collects the remaining documents for review by the University Scholarship Office. The second component is the Conditional Revert protocol, which gives every incomplete file five calendar days to be corrected before a staff member decides its outcome. The third component is committee-side filtering that partitions applicants by school and presents each subcommittee's data under that school's criteria. The fourth is an event-driven notification subsystem that delivers SMS and email alerts via Twilio and SendGrid for submission deadlines, correction requests, and results released by the Office of Admission.

\section{Problem Statement}
\label{sec:problem}

Ateneo de Davao University (AdDU) lacks a digital system to manage the admission process of its internal scholarship tracks. The current reliance on a paper-based procedure creates systemic strain on applicants and on the offices and committees that evaluate them. Specifically, the current process presents the following problems:

\begin{enumerate}
  \item Single-Phase, Paper-Heavy Document Collection: Applicants must assemble and physically submit a complete file at once. A missing or illegible document can lead to an abrupt rejection, with no structured window for correction.
  \item Fragmented Committee Data and Untraceable Decisions: Five School Scholarship Subcommittees apply different selection criteria to approximately 1,326 applicants per cycle, and their data must be collated manually. Decisions are not recorded in a form that traces each status change to the person who made it.
  \item Status and Deadline Communication Deficits: The absence of proactive status updates and deadline reminders leaves applicants unaware of submission windows and correction requests, leading to missed deadlines and incomplete files.
\end{enumerate}

Consequently, this study seeks to answer the following specific research questions:

\begin{enumerate}
  \item What is the perceived reduction in the time and effort AdDU scholarship applicants spend on the application process when using the two-stage digital submission of ScholarPath AdDU compared to the current paper-based process?
  \item How does the ScholarPath AdDU Application Lifecycle workflow perform in terms of (a) functional correctness, specifically whether status transitions, including Conditional Revert and Lapsed, match the expected outputs of a structured Lifecycle Test Case Matrix with no status change lacking a recorded human actor, and (b) perceived usefulness of the Data-Feeding Committee Portal to University Scholarship Office staff and School Scholarship Subcommittee members?
  \item To what extent does the automated notification subsystem reduce missed submission deadlines and uncorrected incomplete files among applicants?
  \item What is the overall system usability and performance efficiency of ScholarPath AdDU as evaluated by users based on the ISO/IEC 25010 functional suitability standards?
\end{enumerate}

\section{Objectives of the Study}
\label{sec:objectives}

General Objective: To design, develop, and evaluate ScholarPath AdDU, a data-feeding scholarship application lifecycle and committee decision-support system for Ateneo de Davao University. The platform aims to replace the paper-based scholarship admission process with a digital workflow that supports, but never replaces, the people who decide.

Specific Objectives:

\begin{enumerate}
  \item To develop a Data-Feeding Committee Portal that partitions applications by school, presents each School Scholarship Subcommittee's data under that school's rule set, and records every status change in an audit log with the acting user and timestamp, such that no module assigns, approves, or disapproves an applicant without a recorded human action.
  \item To implement a two-stage digital submission, consisting of Phase~1 Pre-qualification and Phase~2 Full Verification, into a Document Vault verified entirely within the University Scholarship Office, together with a five-day Conditional Revert protocol for incomplete files.
  \item To implement a proactive notification system that alerts applicants to submission deadlines, correction requests, and released results through automated, event-driven SMS and email delivery.
  \item To administer a Pre-Study Student Problem Validation Survey among current or recent AdDU scholarship applicants to establish baseline data on the perceived time and effort burden of the existing process, and to administer a Post-Prototype Perceived Effort Survey to the same cohort following prototype testing, enabling a direct comparative analysis for Research Questions 1 and 3.
  \item To evaluate the ScholarPath AdDU functional prototype using a task-based usability testing protocol aligned with ISO/IEC 25010 standards, measuring functional suitability, performance efficiency, and system usability via the System Usability Scale (SUS).
\end{enumerate}

\section{Significance of the Study}
\label{sec:significance}

The implementation of ScholarPath AdDU provides practical, structural benefits to key stakeholders within the university. For applicants, the system replaces a single physical submission with a two-stage digital process and gives them a defined window to correct incomplete files instead of facing abrupt rejection. For the University Scholarship Office, the platform digitizes document verification and application tracking, reducing the operational burden of handling approximately 1,326 physical files per cycle. For the five School Scholarship Subcommittees, the platform delivers each school's applicant data already partitioned and filterable under that school's criteria, while leaving every selection decision with the committee. For the Office of Admission, the platform records the handoff of each accepted list and ties applicant notifications to its release of results. Because every status change is logged with the responsible person, the system also supports the traceability expected in institutional quality audits. Additionally, this study contributes to the broader field of higher education management by documenting an architectural framework for migrating scholarship admission from paper to a digital, human-governed workflow under the constraints of national data privacy legislation.

\section{Scope and Limitations}
\label{sec:scope}

This study is bounded by the design, development, and evaluation of a functional web and mobile application prototype for Ateneo de Davao University. The platform's scope is defined by its target users, the scholarship tracks it governs, and the specific administrative functions it digitizes. It does not extend to actual monetary disbursement, integration with the university's finance ledger, or selection for external scholarships. Within these boundaries, the system addresses the digital lifecycle of internal scholarship admission, from applicant registration and two-stage document submission through verification, interview scoring, committee selection, the release of results, and the Academic Trajectory Path of awarded scholars.

Scope. The system scope includes:

\begin{itemize}
  \item Target Audience Coverage: The platform is designed for applicants to AdDU's internal scholarship tracks, awarded scholars, University Scholarship Office staff, interview panels, the five School Scholarship Subcommittees, and the Office of Admission.
  \item Scholarship Coverage: The platform governs the three internal scholarship tracks: the Jubilee Scholarship, Grant-in-Aid (GIA), and Working Scholars. External and government grants are recorded only as a non-selection pathway for disbursement.
  \item Core Functionality:
  \begin{itemize}
    \item A Data-Feeding Committee Portal that partitions applications by school, applies each school's rule set as screening flags and sort orders, and records every human decision in an audit log.
    \item A Document Vault receiving a two-stage digital submission (Phase~1 Pre-qualification and Phase~2 Full Verification), with verification performed entirely within the University Scholarship Office.
    \item A Conditional Revert protocol that gives incomplete files five calendar days for correction, after which a staff member confirms any disapproval.
    \item An automated notification subsystem delivering deadline reminders, correction notices, and released results via SMS and email.
    \item Scholar movement support for inter-track realignment, external grant program transfer, and waitlist reallocation, each confirmed by a person.
  \end{itemize}
  \item Development Approach: The prototype implements a hybrid mobile and web architecture to ensure cross-platform compatibility. The system uses standard web frameworks and cloud-based databases to manage real-time synchronization, cloud storage, and secure user authentication.
  \item Evaluation Scope: Usability testing based on ISO/IEC 25010 standards will involve a purposive sample of approximately 10 to 15 participants drawn from current or recent scholarship applicants, University Scholarship Office staff, and School Scholarship Subcommittee members. The evaluation will use a mixed-method approach, combining task-based observational testing (measuring completion rates and time-on-task) with post-session instruments: the System Usability Scale (SUS), the Post-Prototype Perceived Effort Survey (Appendix~\ref{app:post-survey}) compared against the pre-study baseline (Appendix~\ref{app:pre-survey}), and the Committee Portal Usefulness Survey (Appendix~\ref{app:committee-survey}). Functional correctness will be evaluated separately through a Lifecycle Test Case Matrix.
\end{itemize}

Limitations. This study acknowledges the following realistic constraints that define its boundaries:

\begin{itemize}
  \item Decision-Support Only: The platform structures, filters, and presents applicant data. It does not select, award, approve, or disapprove any applicant. Screening flags and school rule sets inform the University Scholarship Office and the School Scholarship Subcommittees, which retain exclusive authority over every status decision.
  \item Financial Transaction Limitation: The prototype does not handle actual monetary disbursements, bank routing, or direct tuition fee deductions from the university's finance ledger.
  \item External Grant Limitation: External and government grants, such as DOST, local government unit, and private foundation scholarships, are outside the platform's selection process. The platform does not grant external organizations access to the AdDU database, and applicants must still complete any external application through the grantor's own channels.
  \item Manual Verification of Academic Records: The platform does not integrate with the records of the applicant's high school or with the AdDU entrance examination system. The high school general average and entrance exam result are read from uploaded documents and verified manually by University Scholarship Office staff; the accuracy of screening flags therefore depends on that verification.
  \item No Algorithmic Selection: The platform does not implement any algorithm for determining which among multiple qualified applicants should receive a scholarship when applicants exceed the available slots. Sorting and filtering follow criteria set by each subcommittee, and the final selection, including any ranking among equally qualified applicants, remains with the School Scholarship Subcommittee.
  \item User Base Limitation: System access is restricted to applicants to AdDU's internal college scholarship tracks, awarded scholars, and designated university personnel. Basic education units, Graduate School students, and the general public are excluded from this platform.
  \item Testing Scope Limitation: Usability evaluation involves a purposive sample of participants rather than a university-wide deployment. Findings will inform prototype refinement for the capstone requirements but do not constitute definitive statistical claims about the system's efficacy across the entire AdDU population.
  \item Prototype-Based Time Estimation Constraint: The participant time estimates collected through the Post-Prototype Perceived Effort Survey (Appendix~\ref{app:post-survey}) are based on a single, controlled prototype testing session rather than sustained real-world use of the deployed system. Consequently, these estimates may not fully reflect actual time-on-process in practice, as factors such as varying network conditions, individual document readiness, and accumulated platform familiarity over repeated use could produce meaningfully different results from those observed during the evaluation session.
  \item Unconfirmed Procedural Parameters: Several procedural details, including interview panel composition, the starting event of the five-day correction window, and the precedence between the Jubilee grant and program-specific admission minimums, remain under confirmation with the University Scholarship Office. The platform implements these as configurable settings rather than fixed rules.
\end{itemize}
